% ===============================================================
% Section V  --  Evaluation Design
% Owner: Evaluation Lead (integrated)
% ===============================================================
\section{Evaluation Design}
\label{sec:evaluation}

\subsection{Evaluation Setup}
We evaluated three systems, the baseline, a scene-level RAG system, and
an utterance-level RAG system, on a shared set of 15 questions, so that
the contribution of retrieval, and of finer-grained retrieval, could be
compared directly. Evaluation was run with \texttt{evaluate.py}, which
passes every question through all three systems and writes answers and
scores to \texttt{results/evaluation\_results.csv} and a companion JSON
file with the full retrieved chunks. An LLM-as-judge component produced
first-pass scores, which were then reviewed and adjusted manually against
the actual outputs.

\textbf{Configuration note.} The evaluation reported here isolates the
two parts of the pipeline separately. \emph{Retrieval} is independent of
the generation model and is the focus of the retrieval-relevance scores.
For \emph{generation}, the batch harness was run with a lightweight local
model (\texttt{distilgpt2}) as a stand-in generator; the deployable,
demonstrable system described in
Sections~\ref{sec:methodology}--\ref{sec:system} uses Gemma~3~4B. This
distinction matters when reading the results: the retrieval findings
transfer directly to the deployed system, whereas the end-to-end
generation scores reflect the stand-in generator and therefore lower-bound
what the deployed model achieves. Re-running the harness against Gemma~3~4B
is the immediate next step and is discussed in
Section~\ref{sec:results} and Section~\ref{sec:conclusion}.

\subsection{Evaluation Questions}
The 15 questions comprise 5 instructor-provided questions and 10
group-designed questions. The instructor questions are: \emph{Who is
Hamlet?}; \emph{What is the role of Lady Macbeth?}; \emph{What is the
conflict between the Montagues and the Capulets?}; \emph{Why does Macbeth
kill Duncan?}; and \emph{Why does Hamlet delay taking revenge?} The 10
group-designed questions deliberately span a range of question types,
concept explanation, multi-scene tracking, beginner explanation,
cross-character theme, contextual reasoning, and stylised generation, and
were chosen to probe both parts of the system we expected to work well
and parts we expected to struggle with, such as questions spanning
several scenes or requiring metaphorical interpretation. The full set
appears in Appendix~\ref{app:eval-table}.

\subsection{Scoring Criteria}
Each answer was scored from 1 to 5 on each applicable criterion:
\begin{itemize}
  \item \textbf{Correctness} --- is the answer factually right about the
  play?
  \item \textbf{Grounding} --- is it supported by the retrieved passages?
  (Not applicable to the baseline, which retrieves nothing.)
  \item \textbf{Retrieval relevance} --- did the system retrieve the
  right scenes or utterances?
  \item \textbf{Usefulness} --- would a beginner find the answer helpful?
  \item \textbf{Style} --- scored only for the stylised question (Q13).
\end{itemize}
A score of 1 means the answer was wrong or useless; 3 means partially
right with clear gaps; 5 means fully correct, well-grounded, and clearly
useful. The same scale was applied consistently across all questions and
systems.
